\input{../tex-inputs/latex-leseansicht-vorspann}

\section[Arthur, Olga und Louise Schnitzler an Gustav Schwarzkopf, 7. 7. {[}1905{]}]{L04395 Arthur, Olga und Louise Schnitzler an Gustav Schwarzkopf, 7. 7. {[}1905{]}}
\nopagebreak\mylabel{L04395v}
\rehead{ }\normalsize\beginnumbering\briefempfaengerindex{Schwarzkopf, Gustav@\textsc{Schwarzkopf, Gustav}!zzzSchnitzler, Louise@\emph{von Louise Schnitzler}!1905-07-071@{7. 7. {[}1905{]}}|(be}\briefempfaengerindex{Schwarzkopf, Gustav@\textsc{Schwarzkopf, Gustav}!zzzSchnitzler, Olga@\emph{von Olga Schnitzler}!1905-07-071@{7. 7. {[}1905{]}}|(be}\briefempfaengerindex{Schwarzkopf, Gustav@\textsc{Schwarzkopf, Gustav}!zzzSchnitzler, Arthur@\emph{von Arthur Schnitzler}!1905-07-071@{7. 7. {[}1905{]}}|(be}
\toendnotes[C]{\smallbreak\pagebreak[2]}
\correspDesc{Versand  durch Arthur Schnitzler, Olga Schnitzler, Louise Schnitzler am 7. 7. [1905] in Reichenau an der Rax
\newline{}Übermittlung  am 7. 7. [1905] in Reichenau an der Rax
\newline{}Zustellung  am 8. 7. [1905] in Wien 1/1 1
\newline{}Erhalt  durch Gustav Schwarzkopf am 8. 7. [1905] in Tiefer Graben 23}\toendnotes[C]{\smallbreak}
\Standort{DLA, A:Schnitzler, HS.1985.1.1897.}
\physDesc{Bildpostkarte, 243 Zeichen
\newline{}Handschrift Arthur Schnitzler: Bleistift, deutsche Kurrent
\newline{}Handschrift Olga Schnitzler: Bleistift, lateinische Kurrent
\newline{}Handschrift Louise Schnitzler: Bleistift, deutsche Kurrent
\newline{}Versand: 1) Stempel: »\nobreak{}\oindex{Reichenau an der Rax@\textbf{Reichenau an der Rax}|pwk}Reichenau bei Payerbach, 7 7 0\textcolor{gray}{5}, 2{[}–{]}6N\nobreak{}«.   2) Stempel: »\nobreak{}\oindex{I., Innere Stadt@\textbf{I., Innere Stadt}|pwk}Wien 1/1 1, 8/7 {[}05{]}, VIII, Bestellt\nobreak{}«. }\toendnotes[C]{\smallbreak}\pstart{}{\pb}Herrn\pend{}\pstart{}Guſtav Schwarzkopf\pend{}\pstart{}Wien I\oindex{I., Innere Stadt@\textbf{I., Innere Stadt}|pw}\pend{}\pstart{}Tiefer Graben 23\oindex{Wien@\textbf{Wien}!I., Innere Stadt@\textbf{I., Innere Stadt}!Tiefer Graben 23@\textbf{Tiefer Graben 23}|pw}.\pend{}{\bigskip}
\pstart
           \centering{}{\pb}\textcolor{gray}{\textbf{Reichenau\oindex{Reichenau an der Rax@\textbf{Reichenau an der Rax}|pw}
               bei Payerbach, N.-Oe.\oindex{Payerbach@\textbf{Payerbach}|pw}, 500 m Seehöhe.}}\pend
           
\pstart
           \centering{}\textcolor{gray}{\textbf{Kurpark\oindex{Kurpark [Reichenau an der Rax]@\textbf{Kurpark [Reichenau an der Rax]}|pw} mit Teich.}}\pend
           \vspace{1em}
\pstart
           \noindent{}{\pb}Lockt Sie’s nicht? Es{ }ſieht allerdings ganz anders aus,
      aber eigentlich{ }ſchöner. Ein Entschluſs! Wir würden uns ſehr
      freuen.\pend
           
\pstart
           Herzlichſt{\\[\baselineskip]}Ihr{\\[\baselineskip]}\spacefill\mbox{A.}\pend
           \leftskip=0em{}\selectlanguage{ngerman}\vspace{1em}
\pstart
           \noindent{}{[}hs. O. Schnitzler:{]} Herzlichen Gruss!\pend
           \pstart \spacefill\mbox{O. S.}\pend{}\selectlanguage{ngerman}\vspace{1em}
\pstart
           \noindent{}{\pb}{[}hs. L. Schnitzler:{]} Herzliche Grüße von
      \spacefill\mbox{\label{K_L04371-1v}\edtext{Louise Schnitzler}{\lemma{\textnormal{\emph{Louise Schnitzler}}}\Cendnote{\textnormal{Das erlaubt eine
            zuverlässige Datierung der Postkarte über das Datum des Poststempels hinaus, da sie nur am 7. 7. 1905 in Reichenau\oindex{Reichenau an der Rax@\textbf{Reichenau an der Rax}|pwk}
      auf Besuch war; dementsprechend konnte die Karte nicht am Vortag verfasst sein.}}}\label{K_L04371-1}}\pend
           \selectlanguage{ngerman}\endnumbering\briefempfaengerindex{Schwarzkopf, Gustav@\textsc{Schwarzkopf, Gustav}!zzzSchnitzler, Louise@\emph{von Louise Schnitzler}!1905-07-071@{7. 7. {[}1905{]}}|)be}\briefempfaengerindex{Schwarzkopf, Gustav@\textsc{Schwarzkopf, Gustav}!zzzSchnitzler, Olga@\emph{von Olga Schnitzler}!1905-07-071@{7. 7. {[}1905{]}}|)be}\briefempfaengerindex{Schwarzkopf, Gustav@\textsc{Schwarzkopf, Gustav}!zzzSchnitzler, Arthur@\emph{von Arthur Schnitzler}!1905-07-071@{7. 7. {[}1905{]}}|)be}\mylabel{L04395h}
\begin{anhang}
\end{anhang}\newcommand{\dateiname}{L04395}\newcommand{\titel}{Arthur, Olga und Louise Schnitzler an Gustav Schwarzkopf, 7. 7. [1905]}\newcommand{\editorInnen}{Herausgegeben von Jahnke, SelmaMüller, Martin Anton}\input{../tex-inputs/latex-leseansicht-abspann}
